\begin{proof}[Proof of \cref{obj:lem:jms-ace-class-relations}]
Fix the parameter record, an index \(n\), and a model \(m\).  For every code index \(j\), the clipped-code discrepancies
\[
x\longmapsto \bar g_j(x)-g_0(x),
\qquad
x\longmapsto \bar q_j(x)-q_0(x)
\]
are \(P_X\)-measurable, because the supplied codes and the regressions are measurable and clipping preserves measurability.  The covariate marginal \(P_X\) is a probability law, since it is the image of the observed-data probability law under the covariate map. % lean: jms_ace_class_relations

Assume first that \(m\in\mathcal P_{\mathrm{ACE},n}^{\mathrm{JMS}}\).  The positive class-index condition and the structural, range, tail, and cumulant requirements of \(\mathcal P_{\mathrm{NG},n}\) are the corresponding member conditions in \cref{obj:def:jms-ace-class}.  The treatment-code requirement in \cref{obj:def:non-gaussian-class} follows from the JMS ACE treatment radius.  Put
\[
f(x)=\bar g_n(x)-g_0(x).
\]
The JMS ACE radius bound gives a finite \(L^r(P_X)\) bound for \(f\).  Since \(P_X\) is a probability law and \(2\le r\), exponent monotonicity gives \(f\in L^1(P_X)\) and
\[
\|f\|_{L^1(P_X)}
\le
\|f\|_{L^r(P_X)} .
\]
Together with \(\|f\|_{L^r(P_X)}\le\varepsilon_{1,n}\) and the nonnegativity of \(\varepsilon_{1,n}\) at the positive index \(n\), this yields
\[
\int |\bar g_n(x)-g_0(x)|\,dP_X(x)\le \varepsilon_{1,n},
\]
with the displayed integrand integrable.  Thus \(m\in\mathcal P_{\mathrm{NG},n}\). % lean: jms_ace_class_relations

Next assume \(m\in\mathcal P_{\mathrm{NG},n}^{\mathrm{ACE}}\).  The inherited non-Gaussian member conditions supply the positive class-index condition and the common structural, range, tail, and cumulant requirements appearing in \cref{obj:def:jms-ace-class}.  The subclass restrictions in \cref{obj:def:ace-comparison-subclass} give
\[
\|\bar g_n-g_0\|_{L^s(P_X)}\le\varepsilon_{1,n},
\qquad
\|\bar q_n-q_0\|_{L^s(P_X)}\le\varepsilon_{2,n}.
\]
Because \(r\le s\) and \(P_X\) is a probability law, exponent monotonicity gives
\[
\|\bar g_n-g_0\|_{L^r(P_X)}
\le
\|\bar g_n-g_0\|_{L^s(P_X)},
\qquad
\|\bar q_n-q_0\|_{L^r(P_X)}
\le
\|\bar q_n-q_0\|_{L^s(P_X)} .
\]
These are the two JMS ACE code-radius requirements, so
\(m\in\mathcal P_{\mathrm{ACE},n}^{\mathrm{JMS}}\). % lean: jms_ace_class_relations

Finally suppose \(s=r\).  The preceding paragraph gives
\[
m\in\mathcal P_{\mathrm{NG},n}^{\mathrm{ACE}}
\quad\Longrightarrow\quad
m\in\mathcal P_{\mathrm{ACE},n}^{\mathrm{JMS}} .
\]
Conversely, if \(m\in\mathcal P_{\mathrm{ACE},n}^{\mathrm{JMS}}\), the first inclusion gives \(m\in\mathcal P_{\mathrm{NG},n}\), and the two JMS ACE \(L^r(P_X)\) code bounds become exactly the two \(L^s(P_X)\) subclass bounds in \cref{obj:def:ace-comparison-subclass}.  Hence
\[
m\in\mathcal P_{\mathrm{ACE},n}^{\mathrm{JMS}}
\quad\Longleftrightarrow\quad
m\in\mathcal P_{\mathrm{NG},n}^{\mathrm{ACE}} .
\]
% lean: jms_ace_class_relations
\end{proof}
